\chapter{System Design \& Implementation}
\label{ch:implementation}

This chapter describes how the mathematical formulation of Chapter~3 is
realised in code. The guiding requirement, set explicitly by the
supervision process, is that every implemented line be traceable to a
stated equation: the policy-gradient loss in the code \emph{is}
Equation~\ref{eq:reinforce_loss}; the temporal-difference loss
\emph{is} Equation~\ref{eq:dqn_loss}. Both learning algorithms are
implemented from scratch in PyTorch \cite{paszke2019} with the Adam
optimiser \cite{kingma2015adam}; environments follow the Gymnasium API
\cite{brockman2016gym, towers2024gymnasium}. No reinforcement-learning
library performs any part of the final-model training.

\section{Architecture overview}
\label{sec:arch}

The system separates four concerns, each in its own module:

\begin{enumerate}
  \item \textbf{Data pipeline} (\texttt{fetch\_spy\_daily.py}): downloads
        and splits market data into episode arrays; run once, cached to disk.
  \item \textbf{Environment} (\texttt{env\_full.py}): a Gymnasium
        \texttt{gym.Env} implementing the MDP --- state construction, trade
        execution, fees, masking, termination.
  \item \textbf{Agents} (\texttt{reinforce\_trading.py},
        \texttt{dqn\_trading.py}): the two from-scratch learners, exposing
        a common \texttt{policy(obs, mask) $\to$ action} interface after
        training.
  \item \textbf{Evaluation harness} (\texttt{run\_final\_experiments.py}):
        trains every method under an identical budget and evaluates all
        methods and baselines on identical held-out episodes, writing JSON
        results and charts.
\end{enumerate}

Because baselines, learned agents and diagnostics all flow through the same
\texttt{run\_episode} function, no method can be advantaged by a differing
evaluation path. Figure~\ref{fig:pipeline} summarises the flow.

\begin{figure}[h]
  \centering
  \includegraphics[width=\textwidth]{system_pipeline}
  \caption{System architecture: four decoupled modules. The evaluation
  harness scores every learned agent and every baseline through one
  shared episode-runner, so no method can gain from a differing
  evaluation path.}
  \label{fig:pipeline}
\end{figure}

\section{Data pipeline}
\label{sec:data}

Daily adjusted closes for SPY (the S\&P~500 ETF) from January 2018 to
December 2024 are retrieved via the \texttt{yfinance} API. SPY is chosen
deliberately: it is among the most liquid instruments in the world, so
the frictionless-execution assumption of the environment (trades fill at
the closing price) is least violated; and as a broad index fund it
carries no single-company idiosyncrasies. Days are grouped into calendar
months; months with fewer than 15 trading days are dropped. The result is
84 monthly episodes, split \emph{chronologically}
(Figure~\ref{fig:split}): the first 58 months form the training set and
the most recent 26 months (November 2022 -- December 2024) the test set.
The split is never shuffled, so the test period is strictly in the future
of every training sample --- the temporal analogue of avoiding data
leakage (Section~\ref{sec:bg_trading}). Episode arrays and metadata are
cached as \texttt{.npz}/\texttt{.json}, making every experiment exactly
reproducible offline.

\begin{figure}[h]
  \centering
  \includegraphics[width=0.95\textwidth]{chrono_split}
  \caption{Chronological train/test split of the 84 monthly episodes.}
  \label{fig:split}
\end{figure}

The seven-year window is not homogeneous, which is a feature: training
data includes the 2018 volatility spikes, the 2020 COVID crash and
rebound, and the 2022 bear market; the held-out test window covers the
late-2022 trough and the 2023--24 recovery. An agent trained on this
window has seen falling markets --- it cannot blame its test behaviour on
never having met one.

\section{The trading environment}
\label{sec:env_impl}

\texttt{FullStateTradingEnv} subclasses \texttt{gym.Env} with
\texttt{observation\_space} $= \mathbb{R}^9$ and
\texttt{action\_space} $=$ \texttt{Discrete(3)}. Internally it tracks cash
$C_t$, fractional units $h_t$, and the cost basis of the open position
(average-cost accounting), from which the nine state features of
Table~\ref{tab:state} are computed on demand.

\paragraph{Trade execution.} A Buy spends
$\min(C_t,\, \Delta)$ dollars with $\Delta = 0.1\,C_0$, acquiring
$x / (P_t(1+c))$ units; a Sell liquidates $\min(h_t P_t,\, \Delta)$ of
position value, crediting $u P_t (1-c)$. Residual position value below the
\$1 minimum trade size is flushed to cash to avoid dust accumulating
outside the mask logic.

\paragraph{Masking and safety.} \texttt{action\_masks()} returns the legal
action set each step. If a caller nevertheless submits an illegal action,
the environment executes Hold instead and flags the event in \texttt{info};
the evaluation harness accumulates this as the illegal-action diagnostic.

\paragraph{Termination.} At the final bar the environment force-liquidates
any open position at the closing price (fee applied), implementing the
mandatory end-of-episode liquidation of Section~\ref{sec:actions}.

\section{From-scratch REINFORCE}
\label{sec:impl_reinforce}

The policy network is a two-hidden-layer MLP
($9 \to 128 \to 128 \to 3$, $\tanh$ activations). Sampling uses a masked
categorical distribution: illegal logits are set to $-10^9$ before the
softmax, so probability mass is renormalised over legal actions only ---
the network is never trained on, and can never propose, an illegal trade.

One gradient update is performed per episode. The core of the training
loop is four lines, in exact correspondence with the derivation of
Section~\ref{sec:pg_derivation}:

\begin{verbatim}
# rewards-to-go, standardised within the episode
G[t] = r[t] + gamma * G[t+1]
G = (G - G.mean()) / (G.std() + 1e-8)
loss = -(log_probs * G).sum()      # L(theta) = -sum_t log pi(a|s) * G_t
loss.backward(); optimizer.step()  # gradient ascent on J
\end{verbatim}

The returns tensor carries no gradient (it is data, not a function of
$\theta$ in the surrogate loss), matching Step~6 of the derivation.

\section{From-scratch Deep Q-learning}
\label{sec:impl_dqn}

The Q-network is an MLP of the same capacity
($9 \to 128 \to 128 \to 3$, ReLU). The implementation contains the three
standard DQN components, each present for a reason stated in
Section~\ref{sec:dqn}:

\begin{itemize}
  \item \textbf{Replay buffer} (50k transitions, uniform sampling,
        minibatch 64): breaks the temporal correlation of consecutive
        market steps.
  \item \textbf{Target network} (synchronised every 500 steps): stabilises
        the bootstrapped TD target.
  \item \textbf{Masked $\varepsilon$-greedy} ($\varepsilon: 1.0 \to 0.05$
        over the first half of training): exploration restricted to legal
        actions; the $\max$ in the TD target is likewise taken over the
        \emph{stored} legal-action mask of the next state.
\end{itemize}

Storing the next-state mask in the buffer is a subtle but necessary
detail: without it, the target could bootstrap from an action that is
illegal in $s_{t+1}$ (e.g.\ Sell while flat), silently overestimating value.

The core update, in exact correspondence with
Equation~(\ref{eq:dqn_loss}):

\begin{verbatim}
q_next = target_net(s_next)               # Q_phi-minus(s', .)
q_next[~mask_next] = -1e9                 # max over LEGAL actions only
y = r + gamma * (1 - done) * q_next.max(dim=1).values
loss = mse(qnet(s).gather(1, a), y.detach())   # TD loss on Bellman targets
loss.backward(); optimizer.step()
\end{verbatim}

\noindent The \texttt{y.detach()} matters: the target is data, not a
function to differentiate through --- gradients flow only into
$Q_\phi(s,a)$, exactly as the semi-gradient TD rule prescribes
\cite{sutton2018}.

\section{PPO ablation row}
\label{sec:impl_ppo}

The third ablation method, PPO, uses the reference Stable-Baselines3
implementation \cite{raffin2021sb3} (\texttt{MlpPolicy}; learning rate
$3{\times}10^{-4}$,
rollout $n_{\mathrm{steps}}{=}512$, batch 64, 8 epochs, $\gamma{=}0.99$,
$\epsilon{=}0.2$), trained through a thin sampler wrapper
(\texttt{MonthSamplerFullEnv}) whose \texttt{reset()} draws a random
training month --- the same episode distribution the scratch methods see.

One protocol asymmetry is worth stating openly. The scratch methods are
\emph{mask-aware}: illegal logits or Q-values are removed before action
selection during both training and evaluation. SB3's PPO has no native
mask support, so during training the environment's forced-Hold override
handles illegal proposals, and during evaluation an illegal deterministic
action falls back to Hold. This mirrors the pilot-study protocol and, if
anything, mildly disadvantages PPO; the illegal-proposal rate is reported
with the results so the effect is quantified rather than hidden.

\section{Verification and testing}
\label{sec:testing}

Two layers of verification precede any market experiment.

\paragraph{Unit tests.} Eight environment tests
(\texttt{test\_env\_full.py}) assert: the accounting identity
$\sum_t r_t = W_T - W_0$ under random legal play; zero wealth change and
zero trades for do-nothing; the exact fee charge on a flat-price Buy;
forced liquidation leaving zero units; Sell masked while flat with the
illegal-action override engaging; Buy masked once cash is exhausted; state
shape and clock boundary values ($\tau_0 = 0$, $\tau_T = 1$). All pass.

\paragraph{Algorithm validation on known-outcome tasks (Phase 0).} Before
touching market data, the same algorithm implementations were validated on
games where success is unambiguous: CartPole-v1, Flappy Bird and
LunarLander-v3. Every learner clears its random floor (e.g.\ CartPole:
random 28.8 vs.\ REINFORCE 292 and DQN 500/500). A learning algorithm that
cannot solve a game with a known answer cannot be trusted on markets,
where the correct answer is unknown --- this ordering was a fixed
methodological requirement of the project.

\section{Experimental configuration}
\label{sec:config}

\begin{table}[h]
\centering
\caption{Hyperparameters and budgets for the final experiments.}
\label{tab:hparams}
\begin{tabular}{lll}
\toprule
\textbf{Setting} & \textbf{REINFORCE} & \textbf{Deep Q} \\
\midrule
Network & MLP $9{-}128{-}128{-}3$, $\tanh$ & MLP $9{-}128{-}128{-}3$, ReLU \\
Optimiser / lr & Adam, $10^{-3}$ & Adam, $10^{-3}$ \\
Discount $\gamma$ & 0.99 & 0.99 \\
Update cadence & once per episode & every step after warm-up \\
Replay / batch & --- & 50k buffer, batch 64 \\
Target sync & --- & every 500 steps \\
Exploration & sampling from $\pi_\theta$ & $\varepsilon$: $1.0 \to 0.05$ (first 50\%) \\
Budget & 80{,}000 env steps & 80{,}000 env steps \\
Seeds & \{42, 43, 44\} & \{42, 43, 44\} \\
\midrule
Environment & \multicolumn{2}{l}{$C_0 = \$10$k; fee 5 bps; slice $\Delta = 0.1\,C_0$; windows $k{=}5$, $m{=}10$} \\
\bottomrule
\end{tabular}
\end{table}

\section{Reproducibility}
\label{sec:repro}

Every claim in Chapter~5 is regenerable from committed code and cached
data. Specifically: (i) all random number generators (Python, NumPy,
PyTorch) are seeded, and results are reported over three seeds rather
than a favourable one; (ii) the data pipeline caches its output, so
experiments never silently retrain on different data; (iii) each
experiment writes a JSON record containing the configuration, per-month
results, training curves and timings alongside the charts; and (iv) all
experiments run on CPU in minutes (REINFORCE $\approx$ 25\,s, DQN
$\approx$ 70\,s, PPO $\approx$ 35\,s per seed; the entire ablation
including the fee-sensitivity study in $\approx$ 12 minutes) --- a
deliberate property: every result in this dissertation can be regenerated
from scratch during a viva demonstration.

\paragraph{Chapter summary (plain language).} The mathematics of
Chapter~3 now exists as four small, separately testable programs: one
that fetches and organises seven years of real price data; one that
simulates the trading month (the ``game board''), charging fees and
enforcing the rules; and two learners built from the ground up, each of
whose central lines of code is a direct transcription of a derived
equation. Before being allowed near market data, both learners had to
prove themselves on video games with known answers, and the simulator had
to pass eight accounting tests --- including that its per-step rewards
always sum exactly to the month's true profit. The next chapter reports
what happened when they traded.
